\begin{abstract}
Under plane strain, branching initiation from a V-shaped interfacial shear crack
at the corner of a kinked interface between two isotropic elastic materials is
investigated. The crack is treated as a prescribed stationary defect with two
symmetric branches whose faces are in frictionless contact. The singular field
at the crack corner is used as the outer field for a small tensile process zone
oriented along the bisector of one of the materials. Conditions are established
under which compressive contact of the interfacial crack faces and a tensile
normal field in the material containing the process zone are satisfied
simultaneously. Mellin transformation, asymptotic matching, the Wiener--Hopf
method, and the condition of stress regularity at the process-zone tip are used
to derive closed analytical relations for the equilibrium zone length, opening,
work done against separation resistance, and the associated energy parameter.
The dependence of these characteristics on interface angle, elastic contrast,
and load level is examined. Formation of the process zone is shown to weaken
the singularity of the initial configuration containing the V-shaped
interfacial crack, although it does not eliminate it completely. For fixed
geometry, the energy and deformation characteristics of the zone are uniquely
related; hence the corresponding branching-initiation criteria are equivalent
when the critical material parameters are chosen consistently. The results make
it possible to predict initiation of a branched crack once the appropriate
fracture-resistance characteristic of the material is specified independently.
\end{abstract}

\noindent\textbf{Keywords:}
piecewise homogeneous body; kinked interface; V-shaped interfacial shear
crack; process zone; branching initiation; Wiener--Hopf method.
